% GENERATED FILE. Edit canonical lessons, then run npm run generate:books.
% canonical-source-hash: fnv1a64:74284bb562380bbe
% canonical-lessons: PT-C454-remetente, PT-C454-destinatario, PT-C454-portes, PT-C454-expedir, PT-C454-registar

\chapter{Nos correios}
\label{ch:pt-nos-correios}
\hlchaptermodality{454}

\begin{chapteropening}
\textbf{By the end of this chapter:} \emph{I can send things by post: the sender, the addressee, the postage, dispatching and registering a letter.}

Complete the last lesson of chapter 454: I can send things by post: the sender, the addressee, the postage, dispatching and registering a letter.
\end{chapteropening}

\section[o remetente]{o remetente — the sender}

\label{lesson:PT-C454-remetente}



\begin{quote}
Before the new one: say the Portuguese for to haggle, then the Portuguese for to wrap.
\end{quote}

\subsection*{You'll want to know: o remetente}
\textbf{o remetente} — \textquotedblleft{}the sender\textquotedblright{}.

\begin{quote}
O remetente escreve a morada nas costas do envelope.
\end{quote}

\begin{quote}
\emph{The sender writes his address on the back of the envelope.}
\end{quote}

On envelopes: \emph{Rem.} — sender.

\begin{cousinweb}
\emph{remetente} is built on \emph{remeter} 'to send', which goes back to Latin \emph{remittere} 'to send back', on \emph{mittere} 'to send' (Lewis \& Short); EtymDB lists Spanish \emph{remitente} as a cousin.
\end{cousinweb}

\begin{grammarlens}[title={What you've built}]
One more word for telling, explaining and describing at B1.
\end{grammarlens}

\subsection*{Your turn}
\begin{itemize}
\raggedright
  \item \emph{Say it:} \emph{o remetente}
  \item \emph{Say it:} \emph{o remetente}, once more
  \item \emph{Say it:} a sentence of your own with \emph{o remetente}
  \item \emph{Recall:} say the Portuguese for to haggle, then the Portuguese for to wrap, then say \emph{o remetente} again
\end{itemize}

\begin{tcolorbox}[breakable,colback=teal!4,colframe=teal!35!black,title={Before you move on}]
What does o remetente mean? (\textquotedblleft{}The sender\textquotedblright{}.) Say it once more.
\end{tcolorbox}
\section[o destinatário]{o destinatário — the addressee, the recipient}

\label{lesson:PT-C454-destinatario}



\begin{quote}
Before the new one: say the Portuguese for to wrap, then the Portuguese for the sender.
\end{quote}

\subsection*{You'll want to know: o destinatário}
\textbf{o destinatário} — \textquotedblleft{}the addressee, the recipient\textquotedblright{}.

\begin{quote}
A carta voltou porque o destinatário mudou de casa.
\end{quote}

\begin{quote}
\emph{The letter came back because the addressee had moved.}
\end{quote}

\emph{O destino} is where it goes; \emph{o destinatário} is who gets it.

\begin{cousinweb}
\emph{destinatário} is built on \emph{destinar} 'to intend, to send to', a learned borrowing from Latin \emph{dēstināre} 'to make fast, to intend' (Lewis \& Short).
\end{cousinweb}

\begin{grammarlens}[title={What you've built}]
One more word for telling, explaining and describing at B1.
\end{grammarlens}

\subsection*{Your turn}
\begin{itemize}
\raggedright
  \item \emph{Say it:} \emph{o destinatário}
  \item \emph{Say it:} \emph{o destinatário}, once more
  \item \emph{Say it:} a sentence of your own with \emph{o destinatário}
  \item \emph{Recall:} say the Portuguese for to wrap, then the Portuguese for the sender, then say \emph{o destinatário} again
\end{itemize}

\begin{tcolorbox}[breakable,colback=teal!4,colframe=teal!35!black,title={Before you move on}]
What does o destinatário mean? (\textquotedblleft{}The addressee, the recipient\textquotedblright{}.) Say it once more.
\end{tcolorbox}
\section[os portes]{os portes — the postage, the delivery charge}

\label{lesson:PT-C454-portes}



\begin{quote}
Before the new one: say the Portuguese for the sender, then the Portuguese for the addressee.
\end{quote}

\subsection*{You'll want to know: os portes}
\textbf{os portes} — \textquotedblleft{}the postage, the delivery charge\textquotedblright{}.

\begin{quote}
A encomenda custa dez euros mais os portes.
\end{quote}

\begin{quote}
\emph{The parcel costs ten euros plus postage.}
\end{quote}

\emph{Portes grátis} — free delivery.

\begin{cousinweb}
\emph{porte} is built on \emph{portar} 'to carry', from Latin \emph{portāre} 'to carry' (Lewis \& Short): the price of carrying something.
\end{cousinweb}

\begin{grammarlens}[title={What you've built}]
One more word for telling, explaining and describing at B1.
\end{grammarlens}

\subsection*{Your turn}
\begin{itemize}
\raggedright
  \item \emph{Say it:} \emph{os portes}
  \item \emph{Say it:} \emph{os portes}, once more
  \item \emph{Say it:} a sentence of your own with \emph{os portes}
  \item \emph{Recall:} say the Portuguese for the sender, then the Portuguese for the addressee, then say \emph{os portes} again
\end{itemize}

\begin{tcolorbox}[breakable,colback=teal!4,colframe=teal!35!black,title={Before you move on}]
What does os portes mean? (\textquotedblleft{}The postage, the delivery charge\textquotedblright{}.) Say it once more.
\end{tcolorbox}
\section[expedir]{expedir — to dispatch, to send off}

\label{lesson:PT-C454-expedir}



\begin{quote}
Before the new one: say the Portuguese for the addressee, then the Portuguese for the postage.
\end{quote}

\subsection*{You'll want to know: expedir}
\textbf{expedir} — \textquotedblleft{}to dispatch, to send off\textquotedblright{}.

\begin{quote}
A loja expediu a encomenda ontem.
\end{quote}

\begin{quote}
\emph{The shop dispatched the parcel yesterday.}
\end{quote}

\emph{A expedição} — the dispatch, the shipping.

\begin{cousinweb}
\emph{expedir} is a learned borrowing from Latin \emph{expedīre} 'to free the feet, to get ready' (Lewis \& Short), on \emph{pēs} 'foot'.
\end{cousinweb}

\begin{grammarlens}[title={What you've built}]
One more word for telling, explaining and describing at B1.
\end{grammarlens}

\subsection*{Your turn}
\begin{itemize}
\raggedright
  \item \emph{Say it:} \emph{expedir}
  \item \emph{Say it:} \emph{expedir}, once more
  \item \emph{Say it:} a sentence of your own with \emph{expedir}
  \item \emph{Recall:} say the Portuguese for the addressee, then the Portuguese for the postage, then say \emph{expedir} again
\end{itemize}

\begin{tcolorbox}[breakable,colback=teal!4,colframe=teal!35!black,title={Before you move on}]
What does expedir mean? (\textquotedblleft{}To dispatch, to send off\textquotedblright{}.) Say it once more.
\end{tcolorbox}
\section[registar]{registar — to register; to record}

\label{lesson:PT-C454-registar}



\begin{quote}
Before the new one: say the Portuguese for the postage, then the Portuguese for to dispatch.
\end{quote}

\subsection*{You'll want to know: registar}
\textbf{registar} — \textquotedblleft{}to register; to record\textquotedblright{}.

\begin{quote}
Registei a carta para ter um recibo.
\end{quote}

\begin{quote}
\emph{I registered the letter to get a receipt.}
\end{quote}

\emph{Uma carta registada} — a registered letter. In Brazil: \emph{registrar}.

\begin{cousinweb}
\emph{registar} is built on \emph{registo} 'register', from Late Latin \emph{regesta} 'things recorded', on \emph{regerere} 'to carry back, to record' (Lewis \& Short).
\end{cousinweb}

\begin{grammarlens}[title={What you've built}]
One more word for telling, explaining and describing at B1.
\end{grammarlens}

\subsection*{Your turn}
\begin{itemize}
\raggedright
  \item \emph{Say it:} \emph{registar}
  \item \emph{Say it:} \emph{registar}, once more
  \item \emph{Say it:} a sentence of your own with \emph{registar}
  \item \emph{Recall:} say the Portuguese for the postage, then the Portuguese for to dispatch, then say \emph{registar} again
\end{itemize}

\begin{tcolorbox}[breakable,colback=teal!4,colframe=teal!35!black,title={Before you move on}]
What does registar mean? (\textquotedblleft{}To register; to record\textquotedblright{}.) Say it once more.
\end{tcolorbox}
